%%%%%%%%%%%%%%%%%%%%%%%%%%%%%%%%%%%%%%%%%%%%%%%%%%%%%%%%%%%%%%%%%%%%%%
% Lớp 12 | Các số đặc trưng ... ghép nhóm | Khoảng biến thiên, khoảng tứ phân vị của mẫu số liệu ghép nhóm
% Dạng 2D3?1-2: Tìm khoảng biến thiên
% Số câu: 13  (sinh tự động bởi tach_cau.py; từ nay thầy sửa trực tiếp file này)
%%%%%%%%%%%%%%%%%%%%%%%%%%%%%%%%%%%%%%%%%%%%%%%%%%%%%%%%%%%%%%%%%%%%%%

%%%%%%%%%%%%%%%% TRẮC NGHIỆM 4 PHƯƠNG ÁN %%%%%%%%%%%%%%%%

% nguồn: Lop_12_Dai_so_Chuong_3
\begin{ex}%[2D3N1-2]%[ma:8c602256]
	\immini{ 
	Khi thống kê chiều cao (đơn vị là cm) học sinh của lớp $\mathrm{12A}$, người ta  sử dụng mẫu số liệu ghép nhóm và được một kết quả cho bởi bảng bên. Khoảng biến thiên của mẫu số liệu ghép nhóm đó bằng
	 \choice
		{\True  $25$}
		{$5$} 
		 {$20$}
		 {$180$}
		 }
		{\begin{tabular}{|c|c|}
		\hline Nhóm  & Tần số \\
		\hline$[155;160)$  & $2$ \\
		{$[160;165)$}& $5$  \\
		{$[165;170)$}& $21$  \\
		{$[170;175)$}& $11$  \\
		{$[175;180)$}&  $1$ \\
		\hline  & $n=40$ \\
		 \hline
\end{tabular}
} 
\loigiai{ 
	Khoảng biến  thiên $R=180-155=25$.
} 
\end{ex}

% nguồn: Lop_12_Dai_so_Chuong_3
\begin{ex}%[2D3N1-2]%[ma:bc240372]
	Thời gian công tác của $50$ nhân viên một công ty được ghi lại ở bảng số liệu ghép  nhóm sau (đơn vị năm).
	\begin{center}
		\begin{tabular}{|l|c|c|c|c|c|c|}
			\hline Nhóm &{$[2;6)$}&{$[6;10)$}&{$[10;14)$}&{$[14;18)$}&{$[18;22)$}&{$[22;26)$}\\ 
			\hline Tần số & $10$ & $20$ &  $16$ & $3$ & $0$ & $1$\\
			 \hline
		\end{tabular}
	\end{center}
	 Khoảng biến thiên của mẫu số liệu ghép nhóm trên là
	\choice 
		{$18$} 
		 {\True $24$}
		{$14$} 
		 {$21$}
	\loigiai{Khoảng biến thiên của mẫu số liệu ghép nhóm trên  là $26-2=24$.}
\end{ex}

% nguồn: Lop_12_Dai_so_Chuong_3
\begin{ex}%[2D3N1-2]%[ma:a65a085e]
	Bảng sau đây biểu diễn mẫu số liệu ghép nhóm về cân nặng của một số quả bưởi da xanh  ở một khu vực (đơn vị: kg).
	\begin{center}
		\begin{tabular}{|c|c|c|c|c|c|}
			\hline Nhóm & $[1{,}2; 1{,}3)$ &  $[1{,}3; 1{,}4)$ & $[1{,}4; 1{,}5)$ & $[1{,}5; 1{,}6)$ & $[1{,}6; 1{,}7)$ \\
			\hline Tần số & $14$ &  $40$ & $13$ & $10$ & $3$ \\
			 \hline
		\end{tabular}
	\end{center}
	Khoảng biến  thiên của mẫu số liệu ghép nhóm trên là
	 \choice
		 {$0{,}6$}
		{\True $0{,}5$} 
		{$0{,}4$} 
		{$0{,}3$} 
	 \loigiai{
		Khoảng biến thiên của mẫu số liệu ghép  nhóm trên là $R=1{,}7-1{,}2=0{,}5$.
	} 
\end{ex}

% nguồn: Lop_12_Dai_so_Chuong_3
\begin{ex}%[2D3N1-2]%[ma:fcc63038]
	Bảng  sau ghi lại nhiệt độ không khí đo được lúc $12$ giờ trưa hằng ngày tại một trạm quan sát trong $1$ năm.
	\begin{center}
		\begin{tabular}{|l|c|c|c|c|c|}
			 \hline Nhóm &{$[12;16)$}&{$[16;20)$}&{$[20;24)$}&{$[24;28)$}&{$[28;32)$}\\
			\hline Tần số & $85$ &  $63$ & $95$ & $87$ & $35$ \\
			 \hline
		\end{tabular}
	\end{center}
	Khoảng biến thiên của  mẫu số liệu trên là
	\choice 
		{$10$} 
		 {$14$}
		 {$60$}
		{\True  $20$}
	 \loigiai{Khoảng biến thiên của mẫu số liệu trên là $32-12=20$.}
\end{ex}

% nguồn: Lop_12_Dai_so_Chuong_3
\begin{ex}%[2D3N1-2]%[ma:850ad62e]
	Thời gian thực hiện xong một thí nghiệm hoá học của học sinh  lớp $12H$ được ghi lại ở bảng sau:
	\begin{center}
		\begin{tabular}{|c|c|c|c|c|c|}
			\hline \begin{tabular}{c}
			Thời gian \\ 
			 (phút)
		\end{tabular} & {$[5; 6)$} & {$[6; 7)$} & {$[7; 8)$} & {$[8; 9)$} & {$[9; 10)$} \\
		\hline Số học sinh  & $12$ & $25$ & $0$ & $0$ & $1$ \\
		 \hline
	\end{tabular}
\end{center}
Khoảng biến thiên của mẫu số liệu trên  là
 \choice
	{\True $5$} 
	{$4$} 
	{$2$} 
	{$1$} 
\loigiai{ 
	Khoảng biến thiên của mẫu số liệu trên  là $R=10-5=5.$
} 
\end{ex}

% nguồn: Lop_12_Dai_so_Chuong_3
\begin{ex}%[2D3N1-2]%[ma:56dabf5c]
	Một vườn thú ghi lại tuổi thọ (đơn vị: năm) của $20$ con hổ và thu được  kết quả như sau
	\begin{center}
		\begin{tabular}{|c|c|c|c|c|c|}
			 \hline
			Tuổi thọ & $[14; 15)$ & $[15; 16)$ & $[16; 17)$ & $[17; 18)$ & $[18; 19)$ \\ \hline 
			Số con hổ & $1$ & $3$ & $8$ & $6$  & $2$ \\ \hline
		\end{tabular}
	\end{center}
	Khoảng biến thiên của mẫu số liệu ghép nhóm này là 
	 \choice
		{$4$} 
		 {$3$}
		{\True $5$} 
		{$6$} 
	\loigiai{Khoảng biến  thiên của mẫu số liệu ghép nhóm này là $19-14=5$.
	} 
\end{ex}


%%%%%%%%%%%%%%%% ĐÚNG/SAI %%%%%%%%%%%%%%%%

% nguồn: Lop_12_Dai_so_Chuong_3
\begin{ex}%[2D3N1-2]%[2D3H2-2]%[ma:657558ea]
	Bảng dưới đây thống kê độ ẩm không  khí trung bình các tháng năm 2021 tại Đà Lạt và Vũng Tàu (đơn vị: \%).\\
	 \centerline{
	\begin{tabular}{|*{13}{c|}}\hline
		Tháng& $1$& $2$& $3$& $4$& $5$& $6$&  $7$& $8$& $9$& $10$& $11$& $12$\\\hline
		Đà Lạt& $83$& $79$& $79$& $87$& $87$& $87$&  $88$& $89$& $90$& $91$& $88$& $86$\\
		Vũng  Tàu& $75$& $77$& $78$& $77$& $79$& $79$& $81$& $79$& $81$& $83$& $80$& $77$\\\hline
	\end{tabular}
	 }\\
	Người ta lần lượt ghép các số liệu của Đà Lạt, Vũng Tàu thành mẫu số liệu ghép nhóm như  sau:\\
	\centerline{ 
	\begin{tabular}{|c|c|c|c|c|c|}\hline
		\multicolumn{3}{|c|}{Đà lạt}& \multicolumn{3}{c|}{Vũng Tàu}  \\\hline
		Nhóm& Giá trị đại diện& Tần số& Nhóm& Giá trị đại  diện& Tần số\\\hline
		 $[75;78{,}3)$& $76{,}65$& $0$& $[75;78{,}3)$& $76{,}65$& $5$\\
		$[78{,}3;81{,}6)$& $79{,}95$&  $2$& $[78{,}3;81{,}6)$& $79{,}95$& $6$\\
		$[81{,}6;84{,}9)$& $83{,}25$& $1$& $[81{,}6;84{,}9)$& $83{,}25$& $1$\\ 
		$[84{,}9;88{,}2)$& $86{,}55$& $6$& $[84{,}9;88{,}2)$&  $86{,}55$& $0$\\
		$[88{,}2;91{,}5)$& $89{,}85$&  $3$& $[88{,}2;91{,}5)$& $89{,}85$& $0$\\\hline
	\end{tabular}
	 }\\
	Xét tính đúng sai của các khẳng định sau: 
	\choiceTF 
		{\True Xét số liệu ở Đà Lạt ta có  khoảng biến thiên là $16{,}5$}
		{\True Xét số liệu ở Đà Lạt thì ta có độ lệch chuẩn của mẫu số liệu ghép nhóm (làm tròn  kết quả đến hàng phần trăm) là $3{,}28$}
		{Xét số  liệu ở Vũng Tàu thì phương sai của mẫu số liệu ghép nhóm là $4{,}235$}
		{Đà Lạt có nhiệt độ không khí trung  bình tháng đồng đều hơn vì độ lệch chuẩn nhỏ hơn}
	\loigiai{ 
		\begin{itemchoice}
			\itemch Đúng:  Xét số liệu ở Đà Lạt:\\
			Khoảng biến  thiên: $R=91{,}5-75=16{,}5$.
			\itemch Đúng: Số phần tử của  mẫu là $n=12$.\\
			Tần số tích lũy  của các nhóm lần lượt là $cf_1=0,cf_2=2,cf_3=3,cf_4=9,cf_5=12$.\\
			Ta có: $\dfrac{n}{4}=\dfrac{12}{4}=3$ suy ra  nhóm 3 là nhóm đầu tiên có tần số tích lũy lớn hơn hoặc bằng 3. \\
			Xét nhóm 3 là nhóm $[81{,}6;84{,}9)$  có $s=81{,}6$, $h=3{,}3$, $n_3=1$ và nhóm 2 là nhóm $[78{,}3;81{,}6)$ có $cf_2=2$. \\
			Ta có tứ phân vị thứ nhất là: $\mathrm{Q}_1=s+\left(\dfrac{3-cf_2}{n_3}\right)\cdot h =81{,}6+\left(\dfrac{3-2}{1}\right)\cdot  3{,}3=84{,}9$.\\
			Ta có: $\dfrac{3n}{4}=\dfrac{3\cdot 12}{4}=9$ suy ra nhóm 4  là nhóm đầu tiên có tần số tích lũy lớn hơn hoặc bằng 9. \\
			Xét  nhóm 4 là nhóm $[84{,}9;88{,}2)$ có $t=84{,}9$; $l=3{,}3$; $n_4=6$ và nhóm 3 là nhóm $[81{,}6;84{,}9)$ có $cf_3=3$.\\
			Ta có tứ phân vị thứ ba là: $\mathrm{Q}_3=t+\left(\dfrac{9-cf_3}{n_4}\right)\cdot l =84{,}9+\left(\dfrac{9-3}{6}\right)\cdot 3{,}3=88{,}2$.\\ 
			Khoảng tứ phân vị của mẫu số liệu ghép nhóm trên  là: $\mathrm{Q}_3-\mathrm{Q}_1=88{,}2-84{,}9=3{,}3$.\\
			 Số trung bình cộng của mẫu số liệu ghép nhóm là: \\
			 $\overline{x}_1=\dfrac{0\cdot 76{,}65+2\cdot 79{,}95+83{,}25+6\cdot 86{,}55+3\cdot 89{,}85}{12}=86$\\
			Phương  sai của mẫu số liệu ghép nhóm là:\\
			 $s_1^2=\dfrac{0(76{,}65-86)^2+2(79{,}95-86)^2+(83{,}25-86)^2+6(86{,}55-86)^2+3(89{,}95-86)^2}{12}=10{,}7825$.\\
			 Độ lệch chuẩn của mẫu số liệu ghép nhóm là: $s_1=\sqrt{s_1^2}=\sqrt{10{,}7825}\approx 3{,}28$.
			\itemch Xét số liệu ở Vũng Tàu:\\ 
			Khoảng  biến thiên: $R=91{,}5-75=16{,}5$.\\
			 Số phần tử của mẫu là $n=12$.\\
			Tần số tích lũy của các nhóm lần lượt là $cf_1=5$, $cf_2=11$, $cf_3=12$, $cf_4=12$, $cf_5=12$.\\ 
			Ta có: $\dfrac{n}{4}=\dfrac{12}{4}=3$ suy ra nhóm 1 là nhóm đầu tiên có tần số tích lũy lớn hơn hoặc bằng 3. \\ 
			Xét nhóm 1 là nhóm $[75;78{,}3)$ có $s=75,$;  $h=3{,}3$; $n_1=5$.\\
			Ta có tứ phân vị thứ nhất là: $\mathrm{Q}_1=s+\left(\dfrac{3-cf_0}{n_1}\right)\cdot  h =75+\left(\dfrac{3-0}{5}\right)\cdot 3{,}3=76{,}98$.\\
			Ta có: $\dfrac{3n}{4}=\dfrac{3\cdot  12}{4}=9$ suy ra nhóm 2 là nhóm đầu tiên có tần số tích lũy lớn hơn hoặc bằng 9. \\
			Xét nhóm 2 là nhóm $[78{,}3;81{,}6)$ có $t=78{,}3$; $l=3{,}3$; $n_2=6$ và nhóm 1 là nhóm  $[75;78{,}3)$ có $cf_1=5$. \\
			Ta có tứ phân vị thứ ba là: $\mathrm{Q}_3=t+\left(\dfrac{9-cf_1}{n_2}\right)\cdot l =75+\left(\dfrac{9-5}{6}\right)\cdot 3{,}3=77{,}2$.\\ 
			Khoảng  tứ phân vị của mẫu số liệu ghép nhóm trên là: $\mathrm{Q}_3-\mathrm{Q}_1=77{,}2-76{,}98=0{,}22$.\\
			Số trung bình cộng của mẫu số liệu ghép  nhóm là: $\overline{x}_2=\dfrac{5\cdot 76{,}65+6\cdot 79{,}95+83{,}25}{12}=78{,}85$.\\
			 Phương sai của mẫu số liệu ghép nhóm là:\\
			$s_2^2=\dfrac{5(76{,}65-78{,}85)^2+6(79{,}95-78{,}85)^2+(83{,}25-78{,}85)^2}{12}=4{,}235$. 
			\itemch Sai: Độ lệch chuẩn của mẫu số  liệu ghép nhóm là: $s_2=\sqrt{s_2^2}=\sqrt{4{,}235} \approx 2{,}06$.\\
			Vũng Tàu có nhiệt độ không khí trung bình tháng đồng đều  hơn vì độ lệch chuẩn nhỏ hơn.
		\end{itemchoice}
	} 
\end{ex}

% nguồn: Lop_12_Dai_so_Chuong_3
\begin{ex}%[2D3H1-2]%[ma:cca6e78c]
	Bảng sau biểu diễn mẫu số liệu ghép nhóm thống kê mức lương của một công  ty (đơn vị: triệu đồng).
	\begin{center}
		\begin{tabular}{|c|c|}
			\hline Nhóm & Tần  số \\
			\hline $[10; 15)$  & $15$ \\
			\hline  $[15; 20)$ & $18$ \\
			\hline $[20; 25)$ & $10$  \\
			\hline $[25; 30)$  & $10$ \\
			\hline $[30; 35)$  & $5$ \\
			\hline $[35; 40)$  & $2$ \\
			 \hline & $n=60$ \\
			 \hline
		\end{tabular}
	\end{center}
	 Xét tính đúng sai của các khẳng định sau:
	 \choiceTF
		{\True Khoảng biến thiên của mẫu số liệu ghép nhóm trên là $R=30$} 
		{\True Số phần tử  của mẫu là $n=60$}
		{\True Tứ phân vị thứ  nhất là $Q_{1}=15$}
		{Khoảng tứ phân vị của mẫu số liệu ghép nhóm  trên là $\Delta_{Q}=3$}
	 \loigiai{
		\begin{itemchoice}
			\itemch Khoảng biến  thiên của mẫu số liệu ghép nhóm trên là $R=40-10=30$.
			 \itemch Số phần tử của mẫu là $n=60$.
			\itemch Tần số tích lũy của  các nhóm lần lượt là
			\[c f_{1}=15,\,c f_{2}=33,\,c f_{3}=43,\,c  f_{4}=53,\,c f_{5}=58,\,c f_{6}=60.\]
			Ta có $\dfrac{n}{4}=\dfrac{60}{4}=15$, suy ra nhóm $1$ là nhóm đầu tiên có tần số tích lũy lớn hơn hoặc bằng $15$. Xét nhóm $1$ là nhóm $[10; 15]$ có  $s=10$, $h=5$, $n_{1}=15$.\\
			Ta có tứ phân vị thứ  nhất là $Q_{1}=s+\left(\dfrac{15-c f_{0}}{n_{1}}\right) \cdot h=10+\left(\dfrac{15-0}{15}\right) \cdot 5=15$.
			\itemch Ta có $\dfrac{3 n}{4}=\dfrac{3\cdot60}{4}=45$ mà $43<45<53$ suy ra nhóm $4$ là nhóm đầu tiên có tần số tích lũy lớn  hơn hoặc bẳng $45$ . \\
			Xét  nhóm $4$ là nhóm [$25; 30$] có $t=25, \,j=5, \,n_{4}=10$ và nhóm $3$ là nhóm [$20; 25$] có $c f_{3}=43$.\\
			Ta có tứ phân vị  thứ ba là $Q_{3}=t+\left(\dfrac{45-c f_{3}}{n_{4}}\right) \cdot l=25+\left(\dfrac{45-43}{10}\right) \cdot 5=26$.\\
			Khoảng tứ phân vị của mẫu  số liệu ghép nhóm trên là $Q_{3}-Q_{1}=26-15=9$.
	\end{itemchoice}}
\end{ex}

% nguồn: Lop_12_Dai_so_Chuong_3
\begin{ex}%[2D3H1-2]%[ma:99cdb1bf]
	Một công ty giống cây trồng đã thử nghiệm hai phương pháp chăm sóc khác nhau cho cây  hướng dương. Sau hai tuần, người ta thấy cây được chăm sóc theo cả hai phương pháp đều thấp hơn 50 cm.
	\begin{center}
		\begin{tabular}{cc}
			\begin{tikzpicture}[scale=1,>=stealth,font=\footnotesize,line join=round,line cap=round,x=1mm,y=2mm]
					 \def\xmin{0} \def\xmax{60} \def\ymin{0} \def\ymax{15}
					\draw[color=gray!50,help lines,xstep=2mm,ystep=2mm] (\xmin,\ymin)  grid (\xmax,\ymax);
					\draw[->] (\xmin,0)--(\xmax,0) node [below]{cm}; 
					\draw[->] (0,\ymin)--(0,\ymax) node [above]{Tần  số};
					\node at (0,0)  [xshift=-2mm,yshift=-2mm]{$O$};
					\foreach \x in  {10,20,30,40,50}
					\draw (\x,1pt)--(\x,-2pt) node  [below] {$\x$};
					 \foreach \y in {5,10}
					\draw (1pt,\y)--(-2pt,\y) node [left]  {$\y$};
					\fill[blue,opacity=.75]  (0,0)--(0,6)--(10,6)--(10,8)--(20,8)--(20,12)--(30,12)--(30,8)--(40,8)--(40,6)--(50,6)--(50,0);
				\end{tikzpicture}
			 &
			\begin{tikzpicture}[scale=1,>=stealth,font=\footnotesize,line join=round,line cap=round,x=1mm,y=2mm]
					 \def\xmin{0} \def\xmax{60} \def\ymin{0} \def\ymax{15}
					 \draw[color=gray!50,help lines,xstep=2mm,ystep=2mm] (\xmin,\ymin) grid (\xmax,\ymax);
					\draw[->]  (\xmin,0)--(\xmax,0) node [below]{cm};
					 \draw[->] (0,\ymin)--(0,\ymax) node [above]{Tần số};
					 \node at (0,0) [xshift=-2mm,yshift=-2mm]{$O$};
					 \foreach \x in {10,20,30,40,50}
					\draw (\x,1pt)--(\x,-2pt)  node [below] {$\x$};
					\foreach  \y in {5,10}
					\draw (1pt,\y)--(-2pt,\y) node  [left] {$\y$};
					\fill[blue,opacity=.75] (0,0)--(0,13)--(10,13)--(10,6)--(20,6)--(20,2)--(30,2)--(30,6)--(40,6)--(40,13)--(50,13)--(50,0); 
				\end{tikzpicture}
			\\ 
			\begin{tabular}{c}
				Chiều cao của cây chăm sóc\\ theo phương án  A
			\end{tabular}
			 &
			\begin{tabular}{c}
				Chiều cao của cây chăm sóc\\ theo phương  án B
			\end{tabular}
		\end{tabular}
	\end{center}
	 \choiceTF
		{\True Khoảng  biến thiên của chiều cao các cây được chăm sóc theo mỗi phương pháp $A$ và $B$ bằng nhau}
		{\True Trung bình của chiều cao các cây được chăm sóc theo  mỗi phương pháp $A$ và $B$ bằng nhau}
		{Độ lệch chuẩn của  chiều cao các cây được chăm sóc theo phương án $A$ là $12{,}65$(cm)}
		{Dựa vào độ lệch chuẩn thì chiều cao của các loại cây được chăm sóc theo phương án $B$ ít bị  chênh lệch hơn so với phương án $A$}
	\loigiai{ 
		\begin{itemchoice}
			\itemch Đúng: Khoảng biến thiên của chiều cao các  cây được chăm sóc theo mỗi phương pháp $A$ và $B$ bằng nhau và cùng bằng 50.
			\itemch Đúng: Ước tính số trung bình và độ lệch chuẩn của chiều cao các cây được chăm sóc theo mỗi phương pháp. Cỡ mẫu của hai mẫu số liệu thống kê là $N=40$. Ta có bảng tần số ghép nhóm về chiều  cao của cây được chăm sóc theo phương pháp $A$ như sau: \\
			\centerline{ 
			\begin{tabular}{|*{6}{c|}}\hline
				Chiều cao (cm)& $[0;10)$& $[10;20)$& $[20;30)$& $[30;40)$&  $[40;50)$\\\hline
				Giá trị đại diện& $5$& $15$& $25$&  $35$& $45$\\
				Tần  số& $6$& $8$& $12$& $8$& $6$\\\hline
			\end{tabular}
			}\\ 
			Chiều cao trung bình của các cây được chăm  sóc theo phương án $A$ là:\\
			$\overline{x}_A=\dfrac{5\cdot 6+18\cdot  5+25\cdot 12+35\cdot 8+45\cdot 6}{40}=25$.\\
			Ta có bảng tần số ghép nhóm về chiều cao của cây được chăm sóc theo  phương pháp $B$ như sau: \\
			\centerline{ 
			\begin{tabular}{|*{6}{c|}}\hline
				Chiều cao  (cm)& $[0;10)$& $[10;20)$& $[20;30)$& $[30;40)$& $[40;50)$\\\hline
				Giá  trị đại diện& $5$& $15$& $25$& $35$& $45$\\
				Tần  số& $13$& $6$& $2$& $6$& $13$\\\hline
			\end{tabular}
			 }\\
			Chiều cao  trung bình của các cây được chăm sóc theo phương án $B$ là: \\
			 $\overline{x}_B=\dfrac{5\cdot 13+15\cdot 6+25\cdot 2+35\cdot 6+45\cdot 13}{40}=25$ cm.
			\itemch Độ lệch  chuẩn của chiều cao các cây được chăm sóc theo phương án $A$ là:\\
			$s_A=\sqrt{\dfrac{5^2\cdot 6+15^2\cdot 8+25^2\cdot 12+35^2\cdot  8+45^2\cdot 6}{40}-25^2}\approx 12{,}65$.
			\itemch Độ lệch chuẩn của chiều cao các cây được  chăm sóc theo phương án $B$ là:\\
			$s_B=\sqrt{\dfrac{5^2\cdot 13+15^2\cdot 6+25^2\cdot  2+35^2\cdot 6+45^2\cdot 13}{40}-25^2} \approx 17{,}03$ $\rm cm^2$.
		\end{itemchoice}
	} 
\end{ex}

% nguồn: Lop_12_Dai_so_Chuong_3
\begin{ex}%[2D3H1-2]%[ma:8e8cf558]
	Thời gian chờ khám bệnh của hai phòng khám $1$ và phòng khám $2$ ở thành phố X được cho  trong bảng sau:
	\begin{center}
		\begin{tabular}{|c|c|c|c|c|}
			\hline 
			Thời gian  (phút) & $\left[0;5\right)$ & $\left[5;10\right)$ & $\left[10;15\right)$ & $\left[15;20\right)$ \\
			\hline 
			Phòng khám số $1$ \newline (Số bệnh nhân) & $3$  & $12$ & $15$ & $18$ \\
			 \hline
			Phòng  khám số $2$ \newline (Số bệnh nhân) & $5$ & $10 $& $12$ & $0$ \\
			\hline 
		\end{tabular}
	\end{center}
	 Khi đó
	\choiceTF 
		{\True Tổng số bệnh nhân chờ khám bệnh ở phòng khám số $1$ dưới $5$ phút là $3$} 
		{Khoảng  biến thiên của mẫu số liệu ghép nhóm về thời gian chờ khám bệnh của phòng khám số $1$ là $R_1=15$}
		{Khoảng biến thiên của mẫu số liệu ghép nhóm về thời gian  chờ khám bệnh của phòng khám số $2$ là $R_2=20$}
		{Thời gian chờ khám bệnh ở phòng khám số $2$ phân tán hơn thời gian chờ khám bệnh ở phòng khám số $1$} 
	\loigiai{ 
		\begin{itemchoice}
			\itemch 
			Tổng số  bệnh nhân chờ khám bệnh ở phòng khám số $1$ dưới $5$ phút là $3$.
			 \itemch
			Khoảng biến  thiên của mẫu số liệu ghép nhóm về thời gian chờ khám bệnh của phòng khám số $1$ là $R_1=20-0=20$.
			 \itemch
			Khoảng biến thiên của mẫu  số liệu ghép nhóm về thời gian chờ khám bệnh của phòng khám số $2$ là $R_2=15-0=15$.
			 \itemch
			Vì $R_1 > R_2$ nên thời gian khám bệnh ở phòng khám số $1$ phân tán hơn thời gian chờ khám bệnh ở phòng khám  số $2$.
		\end{itemchoice}
	} 
\end{ex}


%%%%%%%%%%%%%%%% TRẢ LỜI NGẮN %%%%%%%%%%%%%%%%

% nguồn: Lop_12_Dai_so_Chuong_3
\begin{bt}%[2D3N1-2]%[ma:e79ab3c1]
	Bảng dưới  biểu diễn mẫu số liệu ghép nhóm về chiều cao (đơn vị: centimét) của $36$ học sinh nam lớp $12$ ở một trường trung học phổ thông.
	\begin{center}
		\begin{tabular}{|c|c|c|c|c|c|}
			 \hline
			Nhóm & $[160;163)$& $[163;166)$& $[166;169)$&  $[169;172)$& $[172;175)$\\
			 \hline
			Tần số & $6$& $11$& $9$&  $7$& $3$\\
			 \hline
		\end{tabular}
	\end{center}
	Tìm khoảng  biến thiên của mẫu số liệu ghép nhóm đó.
	\shortans{15} 
	\loigiai{ 
		Trong mẫu số liệu ghép nhóm đó,  ta có
		\begin{itemize}
			\item Đầu mút trái của nhóm  $1$ là $a_1=160$.
			\item Đầu mút phải của nhóm $5$ là  $a_6=175$.
		\end{itemize}
		Vậy khoảng biến thiên của mẫu  số liệu ghép nhóm đó là
		$R=a_6-a_1=175-160=15$\,(cm). 
	 }
\end{bt}

% nguồn: Lop_12_Dai_so_Chuong_3
\begin{bt}%[2D3H1-2]%[ma:e256ad30]
	Bạn Minh ghi lại thời gian tập  bóng bàn của mình trong $10$ ngày ở bảng sau đây (đơn vị: phút).
	\begin{center}
		\begin{tabular}{|l|l|l|l|l|l|l|l|l|l|}
			\hline $20$ & $21$ & $22$ & $25$ & $28$ & $30$ & $33$ & $35$ & $36$ & $39$  \\
			\hline 
		\end{tabular}
	\end{center}
	Bạn Minh ghép số liệu trên thành $4$ nhóm có độ dài bằng nhau, với nhóm đầu  tiên là $[20; 25)$. Tính hiệu giữa trung bình của mẫu số liệu ghép nhóm và trung bình của mẫu số liệu ban đầu.
	 \shortans{1,1}
	\loigiai{ 
		Trung bình của mẫu số  liệu ban đầu là \[\overline{x}=\dfrac{20+21+22+25+28+30+33+35+36+39}{10}=28{,}9.\]
		\begin{center}
			\begin{tabular}{|l|c|c|c|c|}
				\hline Thời gian tập  & {$[20; 25)$} & {$[25; 30)$} & {$[30; 35)$} & {$[35; 40)$} \\
				\hline Số ngày & 3 &  2 & 2 & 3 \\
				 \hline
			\end{tabular}
		\end{center}
		 Trung bình của mẫu số liệu ghép nhóm là \[\overline{X}=\dfrac{3\cdot 22{,}5+2\cdot 27{,}5+2\cdot 32{,}5+3\cdot 37{,}5}{3+2+2+3}=30.\]
		Hiệu giữa trung bình của mẫu số liệu ghép nhóm và trung  bình của mẫu số liệu ban đầu \[\overline{X}-\overline{x}=30-28{,}9=1{,}1.\]
	 }
\end{bt}

% nguồn: Lop_12_Dai_so_Chuong_3
\begin{bt}%[2D3H1-2]%[ma:5a23c671]
	Chiều cao  của các cây bạch đàn (đơn vị: m) ở một khu vườn được ghi lại ở bång sau.
	\begin{center}
		\begin{tabular}{|l|c|c|c|c|c|}
			\hline Nhóm &{$[8;  9)$}&{$[9; 10)$}&{$[10; 11)$}&{$[11; 12)$}&{$[12; 13)$}\\
			\hline Tần số & $6$  & $9$ & $15$ & $27$ & $3$ \\
			 \hline
		\end{tabular}
	\end{center}
	Tìm khoảng tứ phân vị của mẫu số liệu  ghép nhóm trên (kết quả làm tròn đến hàng phần trăm).
	 \shortans{1,56}
	 \loigiai{
		\begin{center}
			\begin{tabular}{|l|c|c|c|c|c|}
				 \hline Nhóm &{$[8; 9)$}&{$[9; 10)$}&{$[10; 11)$}&{$[11; 12)$}&{$[12; 13)$}\\
				\hline Tần số & $6$  & $9$ & $15$ & $27$ & $3$ \\
				\hline Tần số tích lũy& $6$ & $15$ & $30$ &  $57$ & $60$\\
				\hline 
			\end{tabular}
		\end{center}
		Ta có  $\dfrac{n}{4}=\dfrac{60}{4}=15$.
		\begin{itemize}
			\item Nhóm $2$ là nhóm đầu tiên  có tần số tích lũy lớn hơn hoặc bằng $15$.\\
			Nhóm $2$ có $s=9$, $h=10-9=1$, tần số của nhóm $n_2=9$ và  nhóm $1$ có tần số tích lũy $cf_1=6$. Ta có
			\[Q_1=s+\left(\dfrac{15-cf_1}{n_2}\right)\cdot h=9+\dfrac{15-6}{9}\cdot 1=10.\] 
			\item Nhóm $4$ là nhóm đầu tiên có tần số tích lũy lớn hơn hoặc  bằng $\dfrac{3n}{4}=45$.\\
			Nhóm $4$ có $t=11$, $l=1$, tần số  của nhóm $n_4=27$ và nhóm $3$ có tần số tích lũy là $cf_3=30$. Ta có
			 \[Q_3=t+\left(\dfrac{45-cf_3}{n_4}\right)\cdot l=11+\dfrac{45-30}{27}\cdot 1=\dfrac{104}{9}.\]
		\end{itemize}
		Do đó $\Delta_Q=Q_3-Q_1\approx 1{,}56$. 
	} 
\end{bt}
